\documentclass[12pt]{article}

% --- Packages ---
\usepackage[margin=1in]{geometry} % 1 inch margins
\usepackage{setspace}             % For line spacing
\usepackage{amsmath}              % For math equations
\usepackage{amssymb}              % For math symbols
\usepackage{booktabs}             % For professional tables
\usepackage{graphicx}             % For resizing tables if necessary
\usepackage{float}                % For table positioning
\usepackage{hyperref}             % For hyperlinks
\usepackage{longtable}            % For tables that might break pages
\usepackage{indentfirst}          % Indent first paragraph of sections
\usepackage[utf8]{inputenc}

% --- Document Setup ---
\doublespacing % Double spacing for the main body

\title{\textbf{Predicting Corporate Bond YTM Changes}}
\author{Team members: Yuxi Geng, Yutung Lin, Qinqin Huang}
\date{}

\begin{document}

\maketitle

\section{Objective}

The objective of this project is to predict corporate bond yield to maturity (YTM) changes using firm-level fundamentals, market-based risk factors, and macroeconomic indicators.
By applying various machine learning regression models such as ElasticNet, Boosting, and Multi-layer Perceptron, we aim to improve prediction accuracy compared to traditional linear regression models.
In addition, we used the change direction of YTM movements as a classification target since there are many outliers in the YTM change values that may affect regression model performance. We applied classification models such as Logistic Regression, ElasticNet Classifier, and Boosting Classifier to predict the direction of YTM changes.

\section{Methodology}

\subsection{Data Description}
We utilized data from Wharton Research Data Services (WRDS) to construct our dataset. The data frequency is monthly, and the data range is from July 2002 to February 2025.
We have a total of 50 features from four categories: bond-level features, firm-level features, macroeconomic features, and industry features. The details of the predictors are listed in the appendix.

The target variable is the corporate bond yield to maturity (YTM) change, defined as the difference in YTM over a one-month horizon:
\begin{equation}
\Delta \text{YTM}_i = \text{YTM}_{i,t} - \text{YTM}_{i,t-1}
\end{equation}
where $\text{YTM}_{i,t}$ is the yield to maturity of bond $i$ at the end of month $t$.

For classification, the target variable is defined as:
\begin{equation}
\text{YTM\_Direction}_i = 
\begin{cases}
\text{up}, & \text{if } \Delta \text{YTM}_i > 1e^{-6} \\
\text{neutral}, & \text{if } \Delta \text{YTM}_i \approx 0 \\
\text{down}, & \text{if } \Delta \text{YTM}_i < -1e^{-6}
\end{cases}
\end{equation}

The total number of observations after merging all datasets is approximately 854,769, covering 12,117 unique bonds over the sample period.

The data descriptive statistics of the yield to maturity (YTM) changes and class distribution are as follows:

\begin{table}[H]
\centering
\begin{tabular}{lrrrrrrrr}
\hline
 & \textbf{count} & \textbf{mean} & \textbf{std} & \textbf{min} & \textbf{25\%} & \textbf{50\%} & \textbf{75\%} & \textbf{max} \\
\hline
\textbf{Value} &
854769.00 &
-0.000005 &
0.01444 &
-0.9808 &
-0.00184 &
-0.00003 &
0.00171 &
0.97903 \\
\hline
\end{tabular}
\caption{Descriptive statistics}
\end{table}

\begin{table}[H]
\centering
\begin{tabular}{lrrr}
\hline
 & \textbf{down} & \textbf{up} & \textbf{neutral} \\
\hline
\textbf{Count} & 435398 & 399205 & 20166 \\
\hline
\end{tabular}
\caption{Class distribution}
\end{table}

\subsection{Feature Construction}
After preparing the bond-level, firm-level, sector-level, and macroeconomic datasets, we align all information at a monthly frequency and construct the final panel used for modeling. 

\begin{enumerate}
    \item \textbf{Merging Procedure}
    We combine the datasets using index columns: CUSIP (bond identifier), PERMNO (firm identifier), and Date (end-of-month date).

    \item \textbf{Missing-Value Handling: }
    Missing values are handled based on variable type:
    \begin{itemize}
        \item \textbf{Categorical Variables:} We use cross-sectionally monthly median to impute missing values, preserving the categorical distribution across firms each month. For example: credit rating dummies (AAA\dots D), rating transition indicators, and Compustat data-status flag (costat). 
        \item \textbf{Numeric Variables:} We also use cross-sectionally monthly median to impute missing values for numeric variables, the remaining NaNs are set to be zero.
    \end{itemize}

    \item \textbf{Normalization}
    We apply two layers of normalization.
    First, all bond-level and firm-level variables are rank-normalized within each month. The variables are mapped to [-1, 1] using:
        \begin{equation}
        \text{scaled\_rank} = 2 \times \frac{\text{rank}}{N} - 1
        \end{equation}
        where \text{rank} is the cross-sectional rank of the variable for that month, and $N$ is the total number of non-missing observations that month. 
        Features in this category include stock features, fundamentals, derived ratios, growth variables, and bond-level numeric predictors. 
        For binary variables (e.g., upgrade/downgrade), ranks are assigned such that 1 maps to 1 and 0 maps to $-1$.

    Second, macroeconomic variables are normalized using z-score normalization based on training set statistics in each rolling window to avoid data leakage as well as keep their time-series properties. 
        \begin{equation}
        \text{normalized\_value} = \frac{\text{value} - \mu_{\text{train}}}{\sigma_{\text{train}}}
        \end{equation}
        where $\mu_{\text{train}}$ and $\sigma_{\text{train}}$ are the mean and standard deviation of the macro variable computed from the training set only.
\end{enumerate}

\subsection{Data Splitting}
We used a rolling window forecasting approach. At each year, we used a time-series split to create training, validation, and test sets. The training set consisted of 10 years of data, the validation set and test set each consisted of 1 year of data. At each year, we tuned hyperparameters based on training set and validation set, then forecast target variables in the test set.
When tuning is not involved or cross-validation is adopted, the validation set was combined with the training set to train the final model. For example, to predict YTM changes in 2013, we used data from 2002 to 2011 as the training set, data in 2012 as the validation set, and data in 2013 as the test set. 

\subsection{Modeling Framework}
\subsubsection{Regression Models}
We implemented and compared the following machine learning models for the regression task:
\begin{singlespace}
\begin{table}[H]
\centering
\begin{tabular}{p{4cm} p{9.5cm}}
\toprule
\textbf{Model} & \textbf{Description} \\
\midrule
Linear Regression 
& Baseline model used for comparison. \\
ElasticNet Regression
& Combination of L1 and L2 regularization, with hyperparameters 
$\alpha$ and $l1\_ratio$ tuned via 10-fold cross-validation. \\
Boosting Regressor (LightGBM) 
& Gradient boosting model built from ensembles of weak learners with early stopping. Hyperparameters such as num\_leaves and learning\_rate are tuned using grid search. \\
Stacking Regressor 
& Ensemble of multiple base models, using a linear regression meta-learner. \\
Multi-layer Perceptron 
& Feedforward neural network. \\
& Multi-layer Perceptron & Feedforward neural network with ReLU hidden layers (Adam, MSE loss). \\
\bottomrule
\end{tabular}
\end{table}
\end{singlespace}

\subsubsection{Classification Models}
For classification task, we implemented and compared the following models:

\begin{singlespace}
\begin{table}[H]
\centering
\begin{tabular}{p{4cm} p{9.5cm}}
\toprule
\textbf{Model} & \textbf{Description} \\
\midrule
Logistic Regression & Baseline model used for comparison. \\
ElasticNet Classifier & Combination of L1 and L2 regularization. We tuned alpha and l1\_ratio hyperparameters using 10-fold cross-validation. \\
Boosting Classifier (LightGBM) & Ensemble of weak learners. We tuned num\_leaves and learning\_rate hyperparameters using grid search. \\
Stacking Classifier & Ensemble of multiple models using logistic regression as meta-model \\
\bottomrule
\end{tabular}
\end{table}
\end{singlespace}

\subsection{Model Evaluation}
We used a rolling window evaluation approach and calculated performance metrics on the test set for each year from 2012 to 2024. Then, the metrics were averaged to assess overall model performance.
For regression models, performance metrics included: MSE, MAE, MedAE, and $R^2$
For classification models, performance metrics included: Accuracy, Precision, and Recall

\section{Results and Discussion}

\subsection{Regression Results}
The Elastic Net Regression model does not significantly outperform the Linear Regression benchmark, indicating that regularization alone provides only limited gains in this setting. The LGBM Regressor and the Stacking Regressor achieve slightly lower MSE than the pure linear model, but their average out-of-sample $R^2$ remains close to zero (and sometimes slightly negative), suggesting only modest incremental predictability from these nonlinear ensembles.
The tuned multilayer perceptron (MLP) achieves the highest average out-of-sample $R^2 \approx 0.064$ across the rolling windows, indicating that the nonlinear network is able to capture some additional time-series variation in YTM changes. However, this comes at the cost of slightly higher MSE/MAE compared with the best tree-based and stacked models, and the overall magnitude of the $R^2$ remains modest. This suggests that, given our current feature set and network size, a feed-forward neural network provides only limited improvement in predictive accuracy over simpler linear and ensemble benchmarks.

\begin{table}[h!]
\centering
\begin{tabular}{lcccc}
\hline
Model & MSE & MAE & MedAE & $R^2$ \\
\hline
Linear Regression  & 0.000094 & 0.003226 & 0.002065 &  0.001753 \\
ElasticNet Regression       & 0.000097 & 0.003107 & 0.001900 & -0.019272 \\
Boosting Regressor     & 0.000089 & 0.003069 & 0.001620 & -0.109083 \\
Stacking Regressor & 0.000086 & 0.002959 & 0.001591 & -0.008823 \\
Multilayer Perceptron & 0.000105 & 0.003522 & 0.002168 & 0.064329 \\
\hline
\end{tabular}
\end{table}

\subsection{Classification Results}
In classification tasks, the Boosting Classifier (LightGBM) and Stacking Classifier outperformed the logistic models.

\begin{singlespace}
\begin{table}[H]
\centering
\begin{tabular}{lccc}
\toprule
\textbf{Model} & \textbf{Accuracy} & \textbf{Precision} & \textbf{Recall} \\
\midrule
Logistic Regression & 0.687858 & 0.687190 & 0.687858 \\
ElasticNet Classifier & 0.687915 & 0.686969 & 0.687915 \\
Boosting Classifier (LightGBM) & 0.713470 & 0.719949 & 0.713470 \\
Stacking Classifier & 0.713615 & 0.720544 & 0.713615 \\
\bottomrule
\end{tabular}
\end{table}
\end{singlespace}

\section{Conclusion}
In this project, we explored various machine learning models to predict corporate bond yield to maturity changes using firm-level fundamentals, market-based risk factors, and macroeconomic indicators. Our findings indicate that ensemble methods like LightGBM provide clear gains in the classification task, where they outperform traditional linear models in terms of accuracy, precision, and recall. For the regression task, they also offer modest improvements over linear benchmarks. The MLP model achieves worst MSE suggesting we do not need complex neural networks for this problem given the current feature set.

\appendix
\section{Appendix}

\subsection{Implementation Details}

\begin{enumerate}
    \item \textbf{Software and Libraries}
    \begin{itemize}
        \item CPU: MacBook Pro (Apple M4 Pro)
        \item Programming Language: Python 3.13
        \item Libraries: pandas, numpy, scikit-learn, lightgbm, tensorflow/keras
        \item \href{https://github.com/QinqinAndMacaulayCat/ml_project}{Code Repository (GitHub)}
    \end{itemize}

    \item \textbf{Computation Time For Each Model (Approximate)}

\begin{table}[H]
\centering
\begin{tabular}{lcc}
\toprule
\textbf{Model} & \textbf{Training/Prediction} \\
\midrule
Linear Regression & 3.7 seconds \\
Elastic Net Regression & 9.2 seconds \\
Boosting Regressor & 2 minutes 15 seconds \\
Stacking Regressor & 4 minutes 50 seconds \\
Logistic Regression & 57 seconds \\
Elastic Net Classification & 635 minutess \\
Boosting Classifier & 10 minutes \\
Stacking Classifier & 19 minutes 7.4 seconds \\
Multilayer Perceptron & 38 minutes 48 seconds \\
\bottomrule
\end{tabular}
\end{table}
    \item \textbf{AI Tools Used: } We used ChatGPT for drafting and refining report text, and GitHub Copilot for code suggestions and completions.
\end{enumerate}

\subsection{Predictor List}

\begin{longtable}{l p{10cm}}
\caption{Comprehensive Description of All Predictors}
\label{tab:all_predictors_combined} \\
\toprule
\textbf{Field Name} & \textbf{Description} \\
\midrule
\endfirsthead

\multicolumn{2}{c}{\tablename\ \thetable{} -- Continued from previous page} \\
\toprule
\textbf{Field Name} & \textbf{Description} \\
\midrule
\endhead

\bottomrule
\endfoot

\bottomrule
\endlastfoot

\multicolumn{2}{l}{\textbf{A. Key Identifier Columns}} \\
CUSIP & Committee on Uniform Securities Identification Procedures. Unique 9-digit identifier for each bond. \\
PERMNO & Permanent Number. Unique firm-level identifier. \\
Date & End of Month Date. \\
\addlinespace
\addlinespace

\multicolumn{2}{l}{\textbf{B. Bond-level Features}} \\
\texttt{tmt} & The remaining time to maturity of the bond in years. \\
\texttt{coupon} & The annual coupon rate of the bond. \\
\texttt{t\_spread} & The weighted average bid-ask spread of the bond which measures liquidity. \\
\texttt{return\_eom} & Bond return of last month. \\
Credit Ratings & One-hot encoded variables for ratings including \texttt{AAA}, \texttt{AA}, \texttt{A}, \texttt{BBB}, \texttt{BB}, \texttt{B}, \texttt{CCC}, \texttt{CC}, \texttt{C}, and \texttt{D}. \\
\texttt{upgrade} & Binary variables indicating if the bond was upgraded in the last month. \\
\texttt{downgrade} & Binary variables indicating if the bond was downgraded in the last month. \\
\addlinespace
\addlinespace

\multicolumn{2}{l}{\textbf{C. Firm-level Features: Stock Market}} \\
\texttt{ret} & Monthly stock return accumulated from daily returns. \\
\texttt{vol} & Monthly return volatility, calculated as the standard deviation of daily log returns. \\
\texttt{dvol} & Dollar trading volume, calculated as price multiplied by shares traded. \\
\texttt{turnover} & Share turnover ratio, defined as shares traded divided by shares outstanding. \\
\texttt{bidask} & Bid-ask spread based on daily closing bid and ask quotes, measuring liquidity cost. \\
\texttt{numtrades} & Total number of trades executed in the month. \\
\texttt{mktcap} & Month-end market capitalization (Price $\times$ Shares Outstanding). \\
\texttt{price\_mean} & The average daily closing stock price over the month. \\
\addlinespace
\addlinespace

\multicolumn{2}{l}{\textbf{D. Firm-level Features: Accounting-Based Derived Indicators}} \\
\texttt{log\_atq} & Natural logarithm of total assets, used as a proxy for firm size. \\
\texttt{lev\_total} & Total leverage ratio, calculated as total liabilities divided by total assets. \\
\texttt{equity\_ratio} & Equity-to-assets ratio, calculated as total common equity divided by total assets. \\
\texttt{roa} & Return on Assets (ROA), defined as net income divided by total assets. \\
\texttt{profit\_margin} & Net profit margin, calculated as net income divided by total revenue. \\
\texttt{int\_coverage} & Interest coverage ratio, defined as operating income divided by interest expense. \\
\texttt{mkt\_cap} & Accounting-based market capitalization. \\
\texttt{market\_to\_book} & Market-to-book ratio, calculated as market capitalization divided by book equity. \\
\addlinespace
\addlinespace

\multicolumn{2}{l}{\textbf{E. Firm-level Features: Growth Indicators (Quarter-over-Quarter)}} \\
\texttt{atq\_growth} & Quarter-over-quarter growth rate of total assets. \\
\texttt{revtq\_growth} & Quarter-over-quarter growth rate of total revenue. \\
\texttt{niq\_growth} & Quarter-over-quarter growth rate of net income. \\
\addlinespace
\addlinespace

\multicolumn{2}{l}{\textbf{F. Industry Features}} \\
\texttt{price} & The monthly closing price of the sector ETF corresponding to the firm's Global Industry Classification Standard (GICS) sector. \\
\texttt{return} & The monthly return of the sector ETF, used as a proxy for sector-specific market performance and shocks. \\
\addlinespace
\addlinespace

\multicolumn{2}{l}{\textbf{G. Macroeconomic Features}} \\
\texttt{sp500} & Monthly level of the S\&P 500 index, capturing overall equity market conditions and risk appetite. \\
\texttt{ir3m} & Monthly average 3-month Treasury yield, representing short-term risk-free rates and monetary policy stance. \\
\texttt{ir10y} & Monthly average 10-year Treasury yield, reflecting long-term discount rates relevant for bond valuation. \\
\texttt{vix} & Monthly average VIX index, measuring market-implied volatility and broad risk aversion. \\
\texttt{gdp} & Real GDP level or growth, aligned to each month, summarizing aggregate economic activity. \\
\texttt{cpi} & Consumer Price Index (level or change), measuring inflation conditions relevant for nominal yields and real required returns. \\
\end{longtable}


\end{document}
